\newpage
\appendix


\chapter{Proof-Theoretic Concepts}
\begin{dfn}[Proof System]
    \phantomsection \label{dfn:proof-system}
    A \keyword{proof system} \( \proofsystem{P} \) for a first-order language \( \mathcal{L} \) consists of the following:
    \begin{enumerate}[topsep=2pt,noitemsep,label=\textup{\texttt{(\alph*)}}]
        \item A (possibly infinite) set \( \Lambda^{\proofsystem{P}} \) of logical axioms--- \( \mathcal{L} \)-formulas that are universally true regardless of
            interpretation.
        \item A finite set \( \Phi^{\proofsystem{P}} \) of rules that determine how new \( \mathcal{L} \)-formulas can be derived from an existing
            collection of \( \mathcal{L} \)-formulas. These rules are typically called the \keyword{rules of inference}.
    \end{enumerate}
\end{dfn}

\begin{rem}
    The set \( \Lambda^{\proofsystem{P}} \) is typically defined using a finite collection of \emph{axiom schemas}, which are essentially
    templates representing an entire family of infinitely many axioms.
    For instance, the axiom schema \phantomsection \label{anchor:axiom-schema}
    \( A {\implies} \mleft( B {\implies} A \mright) \) is part of many proof systems.
    For such proof systems, the formulas \( \varphi_{1} {\implies} \mleft( \psi_{1} {\implies} \varphi_{1} \mright) \) and
    \( \varphi_2 {\implies} \mleft( \psi_1 {\implies} \varphi_{2} \mright) \) are both part of \( \Lambda^{\proofsystem{P}} \) provided
    that \( \varphi_1 \), \( \varphi_2 \), \( \psi_1 \), and \( \psi_2 \) are all \( \mathcal{L} \)-formulas.
\end{rem}

\begin{dfn}[Proof]
    \phantomsection \label{dfn:formal-proof}
    Let \( \proofsystem{P} \) be a proof system for a first-order language \( \mathcal{L} \) and \( \varphi \)
    be an \( \mathcal{L} \)-formula. Furthermore, let \( \Gamma \) be a set of \( \mathcal{L} \)-formulas.
    A \keyword{formal proof} of \( \varphi \) from \( \Gamma \) in \( \proofsystem{P} \) is a finite sequence
    \( \sequence{\Sigma_i}{i \in n } \) of length \( n \in \naturalset \) such that the following conditions hold:
    \begin{enumerate}[topsep=2pt,noitemsep,label=\textup{\texttt{(\alph*)}}]
        \item For all \( i \in n \), at least one of the following is true:
            \begin{itemize}[topsep=2pt,noitemsep]
                \item \( \Sigma_i \in \Gamma \)
                \item \( \Sigma_i \in \Lambda^{\proofsystem{P}} \)
                \item \( \Sigma_i \) is derived using one of the rules in \( \Phi^{\proofsystem{P}} \)
            \end{itemize}
        \item \( \Sigma_{n-1} = \varphi \)
    \end{enumerate}
\end{dfn}

\begin{dfn}
    \phantomsection \label{dfn:provability}
    An \( \mathcal{L} \)-formula \( \varphi \) is \keyword{provable} from \( \Gamma \) in \( \proofsystem{P} \) if and only if
    there exists a formal proof of \( \varphi \) from \( \Gamma \) in \( \proofsystem{P} \).
\end{dfn}

\begin{notation}
    \phantomsection \label{notation:provable}
    We write \( \Gamma \derives_{\proofsystem{P}} \varphi \) to mean
    that \( \varphi \) is provable from \( \Gamma \) in \( \proofsystem{P} \).
    The subscript can be omitted if the context makes it clear which proof system is being used.
    We write \( \Gamma \models \varphi \) to mean that \( \varphi \) is a \emph{semantic consequence}
    of \( \Gamma \); that is, for every interpretation where all the \( \mathcal{L} \)-formulas in \( \Gamma \)
    is true, then the \( \mathcal{L} \)-formula \( \varphi \) is also true.
\end{notation}

\begin{dfn}[Soundness]
    \phantomsection \label{dfn:soundness-of-proof-system}
    A proof system \( \proofsystem{P} \) of a first-order language \( \mathcal{L} \) is \keyword{sound}
    if and only if for any set of \( \mathcal{L} \)-formulas \( \Gamma \) and any \( \mathcal{L} \)-formula
    \( \varphi \), \( \Gamma \derives_{\proofsystem{P}} \varphi {\implies} \Gamma \models \varphi \).
\end{dfn}

\begin{dfn}[Completeness]
    \phantomsection \label{dfn:completeness-of-proof-system}
    A proof system \( \proofsystem{P} \) of a first-order language \( \mathcal{L} \) is \keyword{complete}
    if and only if for any set of \( \mathcal{L} \)-formulas \( \Gamma \) and any \( \mathcal{L} \)-formula
    \( \varphi \), \( \Gamma \models \varphi {\implies} \Gamma \derives_{\proofsystem{P}} \varphi \). In other
    words, it is possible to construct a proof in \( \proofsystem{P} \) for every possible logical consequence.
\end{dfn}

\noindent
Most proof systems in practice fall into either one of the following categories:
\begin{enumerate*}[label=(\roman*)]
    \item Hilbert systems,
    \item natural deduction.
\end{enumerate*}
Hilbert systems gravitate towards the use of logical axioms to express the laws of deductive reasoning. The number of rules
of inference in such systems is typically kept at a bare minimum.
In fact, most Hilbert systems include only a single rule of inference---Modus Ponens.
Natural deduction, on the other hand, does the exact opposite by incorporating as many rules of inference as possible in order
to mimic the ``natural'' way of human reasoning.

\vspace{0.3cm}
\noindent
Both Hilbert systems and natural deduction have their merits.
Natural deduction is more human-friendly compared to Hilbert systems.
Rather than basing the reasoning process on a set of seemingly abstruse and esoteric axioms,
natural deduction revolves around a set of rules that model human deductive thought.
That is why proofs written within the framework of natural deduction are often easier to read and understand,
which is a huge boon to mathematicians.
The only downside is that it renders the implementation of proof-checking software far more complicated than it need be,
which is why most proof-checkers are implemented as Hilbert systems.

\vspace{0.3cm}
\noindent
All the proofs in this book are derived via natural deduction.
Refer to Appendix \ref{appendices:natural-deduction} for the full specification of the proof system used in this book.

\chapter{Natural Deduction} \phantomsection \label{appendices:natural-deduction}

\begin{notation}
    Let \( \mathcal{L} \) be any first-order language and \( \proofsystem{P} \) be our proof system for \( \mathcal{L} \).
    The following notational conventions will be used in the specification of \( \proofsystem{P} \):
    \begin{enumerate}[topsep=2pt,noitemsep,label=\textup{\texttt{(\alph*)}}]
        \item Lowercase letters (\( x, y, z\)) denote variables.
        \item Uppercase letters (\( A,B,C,\dots \)) and \( \varphi \) denote \( \mathcal{L} \)-formulas.
        \item Uppercase Greek letters (\( \Delta, \Gamma \)) denote \emph{sets} of \( \mathcal{L} \)-formulas.
        \item \aref{dfn:functional-notation}{Functional notation} is used whenever there is a need to emphasize
            that an \( \mathcal{L} \)-formula has free variables (i.e., \( \varphi(x) \)).
        \item Ellipses are used to indicate that an \( \mathcal{L} \)-formula may contain any number of free variables other
            than the one indicated (i.e., \( \varphi(x,\dots) \)).
        \item We write \( \varphi(x \mapsto t) \) to denote the \( \mathcal{L} \)-formula obtained by replacing every free
            occurrence of \( x \) in \( \varphi \) with \( t \).
    \end{enumerate}
\end{notation}

\subsection*{Logical Axioms}
\( \proofsystem{P} \) has only one logical axiom, and that is the Law of Identity: \( \forall x \, \mleft[ x = x \mright] \).
The rest of \( \proofsystem{P} \) consists of rules for deriving new \( \mathcal{L} \)-formulas from a (possibly empty) set of
existing \( \mathcal{L} \)-formulas.

\subsection*{Rules of Inference}
\begin{enumerate}[topsep=2pt,itemsep=10pt,label=\textup{\texttt{(\arabic*)}}]
    \item {Assumption}: \phantomsection \label{rule-of-inference:assumption}

        If \( A \in \Gamma \), then \( \Gamma \derives A \).

    \item Monotonicity: \phantomsection \label{rule-of-inference:monotonicity}

        If \( \Gamma^{\prime} \derives A \) and \( \Gamma^{\prime} \subset \Gamma \), then \( \Gamma \derives A \).

    \item Substitution: \phantomsection \label{rule-of-inference:substitution}

        If \( \Gamma \derives \mleft( x = y \mright) \), then \( \Gamma \derives \mleft( \varphi(x, \dots) {\iff} \varphi(y, \dots) \mright) \)
        for any formula \( \varphi \).

    \item Double Negation: \phantomsection \label{rule-of-inference:double-negation}

        If \( \Gamma \derives A \), then \( \Gamma \derives \lnot \, \mleft( \lnot \, A \mright) \).

    \item Conjunction Introduction: \phantomsection \label{rule-of-inference:conjunction-introduction}

        If \( \Gamma \derives A \) and \( \Gamma \derives B \), then \( \Gamma \derives \mleft( A \,\,{\land}\,\, B \mright) \).

    \item Conjunction Elimination: \phantomsection \label{rule-of-inference:conjunction-elimination}

        If \( \Gamma \derives \mleft( A \,\,{\land}\,\, B \mright) \), then \( \Gamma \derives A \).

    \item Disjunction Introduction: \phantomsection \label{rule-of-inference:disjunction-introduction}

        If \( \Gamma \derives A \), then \( \Gamma \derives \mleft( A \,\,{\lor}\,\, B \mright) \).

    \item Disjunction Elimination: \phantomsection \label{rule-of-inference:disjunction-elimination}

        If \( \Gamma \derives \mleft( A \,\,{\lor}\,\, B \mright) \), \( \Gamma \derives \mleft( A {\implies} C \mright) \),
        and \( \Gamma \derives \mleft( B {\implies} C \mright) \), then \( \Gamma \derives C \).

    \item Commutation: \phantomsection \label{rule-of-inference:commutation}

        If \( \Gamma \derives \mleft( A \,\,{\land}\,\, B \mright) \), then \( \Gamma \derives \mleft( B \,\,{\land}\,\, A \mright) \).

        \noindent
        If \( \Gamma \derives \mleft( A \,\,{\lor}\,\, B \mright) \), then \( \Gamma \derives \mleft( B \,\,{\lor}\,\, A \mright) \).

    \item Association: \phantomsection \label{rule-of-inference:association}

        If \( \Gamma \derives \mleft( A \,\,{\land}\,\, \mleft( B \,\,{\land}\,\, C \mright) \mright) \), then
        \( \Gamma \derives \mleft( \mleft( A \,\,{\land}\,\, B \mright) \,\,{\land}\,\, C \mright) \).

        \noindent
        If \( \Gamma \derives \mleft( A \,\,{\lor}\,\, \mleft( B \,\,{\lor}\,\, C \mright) \mright) \), then
        \( \Gamma \derives \mleft( \mleft( A \,\,{\lor}\,\, B \mright) \,\,{\lor}\,\, C \mright) \).

        \noindent
        If \( \Gamma \derives \mleft( \mleft( A \,\,{\land}\,\, B \mright) \,\,{\land}\,\, C \mright) \), then
        \( \Gamma \derives \mleft( A \,\,{\land}\,\, \mleft( B \,\,{\land}\,\, C \mright) \mright) \).

        \noindent
        If \( \Gamma \derives \mleft( \mleft( A \,\,{\lor}\,\, B \mright) \,\,{\lor}\,\, C \mright) \), then
        \( \Gamma \derives \mleft( A \,\,{\lor}\,\, \mleft( B \,\,{\lor}\,\, C \mright) \mright) \).

    \item Distribution: \phantomsection \label{rule-of-inference:distribution}

        If \( \Gamma \derives \mleft( A \,\,{\land}\,\, \mleft( B \,\,{\lor}\,\, C \mright) \mright) \), then
        \( \Gamma \derives \mleft( \mleft( A \,\,{\land}\,\, B \mright) \,\,{\lor}\,\, \mleft( A \,\,{\land}\,\, C \mright) \mright) \).

        \noindent
        If \( \Gamma \derives \mleft( A \,\,{\lor}\,\, \mleft( B \,\,{\land}\,\, C \mright) \mright) \), then
        \( \Gamma \derives \mleft( \mleft( A \,\,{\lor}\,\, B \mright) \,\,{\land}\,\, \mleft( A \,\,{\lor}\,\, C \mright) \mright) \).

        \noindent
        If \( \Gamma \derives \mleft( \mleft( A \,\,{\land}\,\, B \mright) \,\,{\lor}\,\, \mleft( A \,\,{\land}\,\, C \mright) \mright) \),
        then \( \Gamma \derives \mleft( A \,\,{\land}\,\, \mleft( B \,\,{\lor}\,\, C \mright) \mright) \).

        \noindent
        If \( \Gamma \derives \mleft( \mleft( A \,\,{\lor}\,\, B \mright) \,\,{\land}\,\, \mleft( A \,\,{\lor}\,\, C \mright) \mright) \),
        then \( \Gamma \derives \mleft( A \,\,{\lor}\,\, \mleft( B \,\,{\land}\,\, C \mright) \mright) \).

    \item Disjunction Definition: \phantomsection \label{rule-of-inference:disjunction-definition}

        If \( \Gamma \derives \mleft( A \,\,{\lor}\,\, B \mright) \), then
        \( \Gamma \derives \mleft( \lnot \, \mleft( \lnot \, A \,\,{\land}\,\, \lnot \, B \mright) \mright) \).

        \noindent
        If \( \Gamma \derives \mleft( \lnot \, \mleft( \lnot \, A \,\,{\land}\,\, \lnot \, B \mright) \mright) \),
        then \( \Gamma \derives \mleft( A \,\,{\lor}\,\, B \mright) \).

    \item Conditional Definition: \phantomsection \label{rule-of-inference:conditional-definition}

        If \( \Gamma \derives \mleft( A {\implies} B \mright) \), then \( \Gamma \derives \mleft( \lnot \, A \,\,{\lor}\,\, B \mright) \).

        \noindent
        If \( \Gamma \derives \mleft( \lnot \, A \,\,{\lor}\,\, B \mright) \), then \( \Gamma \derives \mleft( A {\implies} B \mright) \).

    \item Contraposition: \phantomsection \label{rule-of-inference:contraposition}

        If \( \Gamma \derives \mleft( A {\implies} B \mright) \), then \( \Gamma \derives \mleft( \lnot \, B {\implies} \lnot \, A \mright) \).

        \noindent
        If \( \Gamma \derives \mleft( \lnot \, B {\implies} \lnot \, A \mright) \), then \( \Gamma \derives \mleft( A {\implies} B \mright) \).

    \item Biconditional Introduction: \phantomsection \label{rule-of-inference:biconditional-introduction}

        If \( \Gamma \derives \mleft( A {\implies} B \mright) \) and \( \Gamma \derives \mleft( B {\implies} A \mright) \),
        then \( \Gamma \derives \mleft( A {\iff} B \mright) \).

    \item Biconditional Elimination: \phantomsection \label{rule-of-inference:biconditional-elimination}

        If \( \Gamma \derives \mleft( A {\iff} B \mright) \), then \( \Gamma \derives \mleft( A {\implies} B \mright) \).

        \noindent
        If \( \Gamma \derives \mleft( A {\iff} B \mright) \), then \( \Gamma \derives \mleft( B {\implies} A \mright) \).

    \item Law of Excluded Middle: \phantomsection \label{rule-of-inference:law-of-excluded-middle}

        \( \Gamma \derives \mleft( A \,\,{\lor}\,\, \lnot \, A \mright) \).

    \item De Morgan's Laws: \phantomsection \label{rule-of-inference:de-morgans-laws}

        If \( \Gamma \derives \lnot \, \mleft( A \,\,{\land}\,\, B \mright) \), then
        \( \Gamma \derives \mleft( \lnot \, A \,\,{\lor}\,\, \lnot \, B \mright) \).

        \noindent
        If \( \Gamma \derives \lnot \, \mleft( A \,\,{\lor}\,\, B \mright) \), then
        \( \Gamma \derives \mleft( \lnot \, A \,\,{\land}\,\, \lnot \, B \mright) \).

        \noindent
        If \( \Gamma \derives \mleft( \lnot \, A \,\,{\lor}\,\, \lnot \, B \mright) \), then
        \( \Gamma \derives \lnot \, \mleft( A \,\,{\land}\,\, B \mright) \).

        \noindent
        If \( \Gamma \derives \mleft( \lnot \, A \,\,{\land}\,\, \lnot \, B \mright) \), then
        \( \Gamma \derives \lnot \, \mleft( A \,\,{\lor}\,\, B \mright) \).

    \item Absorption: \phantomsection \label{rule-of-inference:absorption}

        If \( \Gamma \derives \mleft( A {\implies} B \mright) \), then \( \Gamma \derives \mleft( A {\implies} \mleft( A \,\,{\land}\,\, B \mright) \mright) \).

    \item Exportation: \phantomsection \label{rule-of-inference:exportation}

        If \( \Gamma \derives \mleft( \mleft( A \,\,{\land}\,\, B \mright) {\implies} C \mright) \),
        then \( \Gamma \derives \mleft( A {\implies} \mleft( B {\implies} C \mright) \mright) \).

    \item Importation: \phantomsection \label{rule-of-inference:importation}

        If \( \Gamma \derives \mleft( A {\implies} \mleft( B {\implies} C \mright) \mright) \), then
        \( \Gamma \derives \mleft( \mleft( A \,\,{\land}\,\, B \mright) {\implies} C \mright) \).

    \item Modus Ponens: \phantomsection \label{rule-of-inference:modus-ponens}

        If \( \Gamma \derives \mleft( A {\implies} B \mright) \) and \( \Gamma \derives A \), then \( \Gamma \derives B \).

    \item Modus Tollens: \phantomsection \label{rule-of-inference:modus-tollens}

        If \( \Gamma \derives \mleft( A {\implies} B \mright) \) and \( \Gamma \derives \lnot \, B \), then \( \Gamma \derives \lnot \, A \).

    \item Disjunctive Syllogism: \phantomsection \label{rule-of-inference:disjunctive-syllogism}

        If \( \Gamma \derives \mleft( A \,\,{\lor}\,\, B \mright) \) and \( \Gamma \derives \lnot \, B \), then \( \Gamma \derives A \).

    \item Hypothetical Syllogism: \phantomsection \label{rule-of-inference:hypothetical-syllogism}

        If \( \Gamma \derives \mleft( A {\implies} B \mright) \) and \( \Gamma \derives \mleft( B {\implies} C \mright) \), then
        \( \Gamma \derives \mleft( A {\implies} C \mright) \).

    \item Constructive Dilemma: \phantomsection \label{rule-of-inference:constructive-dilemma}

        If \( \Gamma \derives \mleft( A {\implies} C \mright) \), \( \Gamma \derives \mleft( B {\implies} D \mright) \), and
        \( \Gamma \derives \mleft( A \,\,{\lor}\,\, B \mright) \), then \( \Gamma \derives \mleft( C \,\,{\lor}\,\, D \mright) \).

    \item Deduction Theorem: \phantomsection \label{rule-of-inference:deduction-theorem}

        If \( \Gamma \union \mleft\{ A \mright\} \derives B \), then \( \Gamma \derives \mleft( A {\implies} B \mright) \).

    \item Reductio ad Absurdum: \phantomsection \label{rule-of-inference:reductio-ad-absurdum}

        If \( \Gamma \union \mleft\{ A \mright\} \derives B \) and \( \Gamma \union \mleft\{ A \mright\} \derives \lnot \, B \), then
        \( \Gamma \derives \lnot \, A \).

    \item Universal Generalization: \phantomsection \label{rule-of-inference:universal-generalization}

        If \( \Gamma \derives \varphi(t, \dots) \) and \( t \) does not appear in any of the formulas in \( \Gamma \),
        then \( \Gamma \derives \forall x \, \varphi(t \mapsto x, \dots) \) provided that \( x \) does not appear in
        \( \varphi(t, \dots) \).

    \item Universal Instantiation: \phantomsection \label{rule-of-inference:universal-instantiation}

        If \( \Gamma \derives \forall x \, \varphi(x,\dots) \) and \( t \) is not a bound variable in \( \varphi(x, \dots) \),
        then \( \Gamma \derives \varphi(x \mapsto t,\dots) \).

    \item Existential Generalization: \phantomsection \label{rule-of-inference:existential-generalization}

        If \( \Gamma \derives \varphi(t,\dots) \) and \( x \) does not appear in \( \varphi(t, \dots) \),
        then \( \Gamma \derives \exists x \, \varphi(t \mapsto x, \dots )\).

    \item Existential Instantiation: \phantomsection \label{rule-of-inference:existential-instantiation}

        If \( \Gamma \derives \exists x \, \varphi(x, \dots) \) and \( t \) does not appear in \( \varphi(x, \dots) \),
        then \( \Gamma \derives \varphi(t \mapsto x, \dots) \).

    \item Existential Definition: \phantomsection \label{rule-of-inference:existential-definition}

        If \( \Gamma \derives \exists x \, \varphi(x, \dots) \), then \( \Gamma \derives \lnot \, \forall x \, \lnot \, \varphi(x, \dots) \).

        \noindent
        If \( \Gamma \derives \lnot \, \forall x \, \lnot \, \varphi(x, \dots) \), then \( \Gamma \derives \exists x \, \varphi(x, \dots) \).
\end{enumerate}

\vspace{0.3cm}
\noindent
The validity of rules \crefrange{rule-of-inference:assumption}{rule-of-inference:reductio-ad-absurdum} is evident from basic propositional logic
(i.e., to verify them one only has to look at the truth tables of the relevant logical operators).
On the other hand, the restrictions imposed on the rules involving quantifiers warrant further discussion.

\vspace{0.3cm}
\noindent
First, we have the \aref{rule-of-inference:universal-generalization}{Universal Generalization} rule, upon which the
following restrictions are imposed:
\begin{enumerate}[label=\textup{\texttt{(\alph*)}}]
    \item \( t \) is not a variable in any of the formulas in \( \Gamma \). \phantomsection \label{enum:universal-generalization-restriction1}
    \item \( x \) is not a variable in \( \varphi \). \phantomsection \label{enum:universal-generalization-restriction2}
\end{enumerate}
Both of these restrictions are necessary to ensure the \aref{dfn:soundness-of-proof-system}{soundness} of \( \proofsystem{P} \).
To illustrate why, suppose that \( \Gamma = \mleft\{ \varphi, \psi, \phi \mright\} \) is a set of \( \mathcal{L} \)-formulas with the formulas
\( \varphi \), \( \psi \), and \( \phi \) defined as follows:
\begin{align*}
    \varphi(t) &\coloneq \text{\( t \) is a cat}\\
    \psi(t,x) &\coloneq \text{\( t \) is an offspring of \( x \)}\\
    \phi &\coloneq \text{\( \forall x \, \exists t \, \psi(x,t) \)}
\end{align*}
Since \( \varphi(t) \in \Gamma \), we can validly infer by \aref{rule-of-inference:assumption}{Assumption} that \( \Gamma \derives \varphi(t) \).
However, without restriction \ref{enum:universal-generalization-restriction1}, we are then permitted to fallaciously apply
\aref{rule-of-inference:universal-generalization}{Universal Generalization} to conclude that
\( \Gamma \derives \forall x \, \varphi(x) \), which is tantamount to the following argument:
\begin{equation*}
    \begin{aligned}
        \text{\bfseries Premise:}&\\
        \text{\bfseries Conclusion:}&
    \end{aligned}
    \begin{aligned}
        &\quad\text{\( t \) is a cat.}\\
        &\quad\text{Therefore, every animal is a cat.}
    \end{aligned}
\end{equation*}
\noindent
Absurd! Just because some animal named \( t \) is a cat does not necessarily mean that every animal in existence is a cat.
In order to validly infer \( \forall z \, \varphi(z) \) from \( \Gamma \),
one must be able to derive \( \varphi \) from \( \Gamma \) using \emph{any}
variable symbol from the language of \( \mathcal{L} \), not just \( t \).
It is easy to see how restriction \ref{enum:universal-generalization-restriction1} enforces this rule.
Since the variable \( t \) is a free variable in at least one formula in \( \Gamma \),
the \aref{rule-of-inference:universal-generalization}{Universal Generalization} rule is not applicable,
so one cannot infer \( \forall z \, \varphi(z) \) from \( \varphi(t) \). Instead, one has to
derive \( \varphi \) from \( \Gamma \) not by mere \aref{rule-of-inference:assumption}{Assumption}, but
by showing that \( \varphi \) still holds even after replacing \( t \) with an arbitrary variable in the language of \( \mathcal{L} \).

\vspace{0.3cm}
\noindent
Restriction \ref{enum:universal-generalization-restriction2} is equally important for the soundness of \( \proofsystem{P} \).
If restriction \ref{enum:universal-generalization-restriction1}'s role is to prevent erroneous inferences, then restriction
\ref{enum:universal-generalization-restriction2}'s role is to preserve the \emph{meaning} of the original formula during the
generalization. To understand this, suppose that \( \Gamma \derives \psi(a,x) \). Since \( a \) is not mentioned in any of the
formulas in \( \Gamma \), restriction \ref{enum:universal-generalization-restriction1} is not violated, so one concludes (albeit fallaciously)
by \aref{rule-of-inference:universal-generalization}{Universal Generalization} that \( \forall x \, \psi(a \mapsto x, x) \).
Informally, this problematic inference translates to:
\begin{equation*}
    \begin{aligned}
        \text{\bfseries Premise:}&\\
        \text{\bfseries Conclusion:}&
    \end{aligned}
    \begin{aligned}
        &\quad \text{\( a \) is an offspring of \( x \).}\\
        &\quad\text{Therefore, every \( x \) is an offspring of \( x \).}
    \end{aligned}
\end{equation*}
\noindent
Again, absurd!
The fact that an animal \( a \) is an offspring of some other animal \( x \) does not logically imply that every animal is an offspring of itself.
There is simply no logical connection whatsoever between the conclusion and the premise. So what exactly went wrong?

\vspace{0.3cm}
\noindent
The problem here is that the variable \( x \) is already a free variable of \( \psi \) before the universal quantifier is applied,
and whenever we replace a free variable in a formula with another free variable in the same formula, we alter the meaning of the formula entirely.
\begin{equation*}
    \begin{aligned}
        \psi(a, x):&\\
        \psi(a \mapsto x, x):&
    \end{aligned}
    \begin{aligned}
        &\quad\text{\( a \) is an offspring of \( x \).}\\
        &\quad\text{\( x \) is an offspring of \( x \).}
    \end{aligned}
\end{equation*}
\noindent
Needless to say, \( \psi(a,x) \) and \( \psi(a \mapsto x, x) \) have completely different meanings.
\( \psi(a, x) \) asserts that the first animal \( a \) is an offspring of
the second animal \( x \), whereas \( \psi(a \mapsto x, x) \) asserts that an animal is an offspring of itself;
and when the universal quantifier is applied to \( \psi(a \mapsto x, x) \), we end up with a formula that says
\emph{every animal is an offspring of itself} instead of the intended
\emph{every animal is an offspring of \( x \)}.

\vspace{0.3cm}
\noindent
Of course, restriction \ref{enum:universal-generalization-restriction2} resolves this by prohibiting the reuse of variable names in the same formula;
that is, we are no longer allowed to infer \( \forall x \, \psi(a \mapsto x, x) \) from \( \psi(a, x) \)
because the symbol \( x \) is already a free variable in \( \psi(a, x) \).
Instead, we must use something other than \( x \) for the bound variable---i.e., \( \forall y \, \psi(a \mapsto y, x) \).

\vspace{0.3cm}
\noindent
To understand the restriction on \aref{rule-of-inference:universal-instantiation}{Universal Instantiation}, consider
the invalid inference from \( \forall x \, \exists t \, \psi(x,t) \) to \( \exists t \, \psi(x \mapsto t, t) \),
which results in the following nonsensical argument:
\begin{equation*}
    \begin{aligned}
        \text{\bfseries Premise}&:\\
        \text{\bfseries Conclusion}&:
    \end{aligned}
    \begin{aligned}
        &\quad\text{Every animal has a parent.}\\
        &\quad\text{There exists an animal that is its own parent.}
    \end{aligned}
\end{equation*}
Since \( t \) is a bound variable in \( \exists t \, \psi(x,t) \)
and the free variable \( x \) appears within the scope of the existential quantifier to which \( t \) is bound,
replacing all occurrences of \( x \) with \( t \) will alter the meaning of the formula because the variable that used
to be \( x \) will no longer be a free variable after the substitution.
\begin{equation*}
    \begin{aligned}
        \exists t \, \psi(x,t)&:\\
        \exists t \, \psi(x \mapsto t, t)&:
    \end{aligned}
    \begin{aligned}
        &\quad \exists t \, \big[ \text{\( x \) is an offspring of \( t \)} \big]\\
        &\quad \exists t \, \big[ \text{\( t \) is an offspring of \( t \)} \big]
    \end{aligned}
\end{equation*}
The variable that used to be \( x \) is supposed to still be a free variable after the substitution; but because
the new variable is called \( t \) and \( t \) is already bound to an existential quantifier, the resulting variable
becomes a bound variable too. A variable with a different name has to be used to retain the intended meaning of the formula:
\begin{equation*}
    \begin{aligned}
        \text{\bfseries Premise}&:\\
        \text{\bfseries Conclusion}&:
    \end{aligned}
    \begin{aligned}
        &\quad \forall x \, \exists t \, \psi(x,t)\\
        &\quad \exists t \, \psi(y, t)
    \end{aligned}
\end{equation*}
Note that the use of the variable \( t \) for \aref{rule-of-inference:universal-instantiation}{Universal Instantiation} is still allowed if \( t \) is a free variable
in the original formula. For example, if \( \Gamma \derives \forall x \, \psi(x, t) \), then \( \Gamma \derives \psi(t, t) \) is a valid inference because \( t \)
is a free variable in \( \psi(x,t) \) and not a bound variable. This is in contrast with restriction \ref{enum:universal-generalization-restriction2} for
\aref{rule-of-inference:universal-generalization}{Universal Generalization} (take note!)
which prohibits the reuse of any variable, free or bound, from the original formula.

\chapter{Tautologies}
\begin{thm}
    \label{tautology:grouping-disjuncts-into-universally-quantified-formula}
    Given any well-formed formulas \( \Gamma \) and \( \varphi \)
    with free variables \( \alpha_{1}, \dots, \alpha_{n}\) and
    \(\beta_{1}, \dots, \beta_{m} \) respectively,
    \begin{equation*}
        \Gamma(\alpha_{1}, \dots, \alpha_{n}) \lor
        \forall x \big( \varphi(x, \beta_{1}, \dots, \beta_{m}) \big)
        \equiv
        \forall x \big( \Gamma(\alpha_{1}, \dots, \alpha_{n}) \lor
        \varphi(x, \beta_{1}, \dots, \beta_{m})  \big)
    \end{equation*}
\end{thm}

\begin{proof}
    Let \( \Gamma \) and \( \varphi \) be well-formed formulas interpreted in the context of an arbitrary, non-empty domain \( D \).
    Let \( \alpha_{1}, \dots, \alpha_{n} \) and \( \beta_{1}, \dots, \beta_{m} \) represent objects
    in \( D \).

    \vspace{0.3cm}
    \noindent Suppose that the formula,
    \begin{equation}
        \label{eq:lhs-statement}
        \Gamma(\alpha_{1}, \dots, \alpha_{n}) \lor
        \forall x \big( \varphi(x, \beta_{1}, \dots, \beta_{m}) \big)
    \end{equation}
    holds. From the truth table for the logical-\( \lor \) operator,
    it is clear that there are only three possible cases in which
    \eqref{eq:lhs-statement} is true.

    % Case 1
    \vspace{0.5cm}
    \noindent {\scshape Case 1}: \( \Gamma(\alpha_{1}, \dots, \alpha_{n}) \land
    \neg \, \forall x \big( \varphi(x, \beta_{1}, \dots, \beta_{m}) \big) \).

    \vspace{0.1cm}
    \noindent
    Since \( \Gamma(\alpha_{1}, \dots, \alpha_{n}) \) holds, then by Disjunction Introduction,
    \begin{equation*}
        \Gamma(\alpha_{1}, \dots, \alpha_{n}) \lor
        \varphi(x_{0}, \beta_{1}, \dots, \beta_{m})
    \end{equation*}
    for any \( x_{0} \in D \). Since \( x_{0} \) is arbitrary, \( \forall x \big( \Gamma(\alpha_{1}, \dots, \alpha_{n})
    \lor \varphi(x, \beta_{1}, \dots, \beta_{m})\big) \) follows from Universal Generalization.


    % Case 2
    \vspace{0.5cm}
    \noindent {\scshape Case 2}: \( \neg \, \Gamma(\alpha_{1}, \dots, \alpha_{n}) \land
    \forall x \big( \varphi(x, \beta_{1}, \dots, \beta_{m}) \big)\).

    \vspace{0.1cm}
    \noindent
    Let \( x_{0} \in D \) be arbitrary. Since \( \forall x \big( \varphi(x, \beta_{1}, \dots, \beta_{m}) \big) \)
    is true, then \( \varphi(x_{0}, \beta_{1}, \dots, \beta_{m}) \) follows from Universal Instantiation. By Disjunction
    Introduction, we thus have,
    \begin{equation*}
        \Gamma(\alpha_{1}, \dots, \alpha_{n}) \lor \varphi(x_{0}, \beta_{1}, \dots, \beta_{m})
    \end{equation*}
    Since \( x_{0} \) is arbitrary, we have \( \forall x \big( \Gamma(\alpha_{1}, \dots, \alpha_{n})
    \lor
    \varphi(x, \beta_{1}, \dots, \beta_{m})\big) \) as needed.


    % Case 3
    \vspace{0.5cm}
    \noindent {\scshape Case 3}: \( \Gamma(\alpha_{1}, \dots, \alpha_{n}) \land
    \forall x \big( \varphi(x, \beta_{1}, \dots, \beta_{m}) \big)\)

    \vspace{0.1cm}
    \noindent
    Let \( x_{0} \in D \) be arbitrary. Clearly, \( \varphi(x_{0}, \beta_{1}, \dots, \beta_{m}) \) is true. Thus,
    \begin{equation*}
        \Gamma(\alpha_{1}, \dots, \alpha_{n}) \lor \varphi(x_{0}, \beta_{1}, \dots, \beta_{m})
    \end{equation*}
    follows from Disjunction Introduction. With \( x_{0} \) being arbitrary, we conclude by Universal Generalization
    that \( \forall x \big( \Gamma(\alpha_{1}, \dots, \alpha_{n}) \lor \varphi(x, \beta_{1}, \dots, \beta_{m}) \big) \).


    \vspace{0.5cm}
    \noindent
    In all three cases, \( \forall x \big( \Gamma(\alpha_{1}, \dots, \alpha_{n}) \lor
    \varphi(x, \beta_{1}, \dots, \beta_{m})\big) \) follows from the initial assumption that
    \( \Gamma(\alpha_{1}, \dots, \alpha_{n}) \lor \forall x \big( \varphi(x, \beta_{1}, \dots, \beta_{m}) \big) \).
    \begin{equation*}
        \therefore \,
        \Gamma(\alpha_{1}, \dots, \alpha_{n})
        \lor
        \forall x \big( \varphi(x, \beta_{1}, \dots, \beta_{m}) \big)
        \implies
        \forall x \big( \Gamma(\alpha_{1}, \dots, \alpha_{n}) \lor
        \varphi(x, \beta_{1}, \dots, \beta_{m})\big)
    \end{equation*}
    Furthermore, this implication holds regardless of any interpretation.
    Thus,
    \begin{equation}
        \label{eq:disjunct-into-universal-quantifier-forward-entailment}
        \Gamma(\alpha_{1}, \dots, \alpha_{n}) \lor
        \forall x \big( \varphi(x, \beta_{1}, \dots, \beta_{m}) \big)
        \models
        \forall x \big( \Gamma(\alpha_{1}, \dots, \alpha_{n})
        \lor
        \varphi(x, \beta_{1}, \dots, \beta_{m})\big)
    \end{equation}

    \noindent
    Next, we show that the logical consequence also holds in the opposite direction. For that, we suppose that
    the formula,
    \begin{equation}
        \label{eq:rhs-statement}
        \forall x \big( \Gamma(\alpha_{1}, \dots, \alpha_{n}) \lor
        \varphi(x, \beta_{1}, \dots, \beta_{m})\big)
    \end{equation}
    holds. Clearly, there are only 2 possible cases in which \eqref{eq:rhs-statement} can be true.

    \vspace{0.5cm}
    \noindent
    {\scshape Case 1}: \( \Gamma(\alpha_{1}, \dots, \alpha_{n}) \) holds.

    \vspace{0.1cm}
    \noindent
    {\scshape Case 2}: \( \neg \, \Gamma(\alpha_{1}, \dots, \alpha_{n}) \land
    \forall x \big( \varphi(x, \beta_{1}, \dots, \beta_{m}) \big)\).

    \vspace{0.5cm}
    \noindent
    In both cases,
    \begin{equation*}
        \Gamma(\alpha_{1}, \dots, \alpha_{n})
        \lor
        \forall x \big( \varphi(x, \beta_{1}, \dots, \beta_{m}) \big)
    \end{equation*}
    follows trivially by Disjunction Introduction.
    \begin{equation*}
        \therefore \,
        \forall x \big( \Gamma(\alpha_{1}, \dots, \alpha_{n})
        \lor
        \varphi(x, \beta_{1}, \dots, \beta_{m}) \big)
        \implies
        \Gamma(\alpha_{1}, \dots, \alpha_{n})
        \lor
        \forall x \big( \varphi(x, \beta_{1}, \dots, \beta_{m}) \big)
    \end{equation*}
    Since the implication is independent of the interpretation, we have,
    \begin{equation}
        \label{eq:disjunct-into-universal-quantifier-backward-entailment}
        \forall x \big( \Gamma(\alpha_{1}, \dots, \alpha_{n}) \lor
        \varphi(x, \beta_{1}, \dots, \beta_{m})\big)
        \models
        \Gamma(\alpha_{1}, \dots, \alpha_{n})
        \lor
        \forall x \big( \varphi(x, \beta_{1}, \dots, \beta_{m}) \big)
    \end{equation}

    \vspace{0.5cm}
    \noindent
    The theorem follows directly from \eqref{eq:disjunct-into-universal-quantifier-forward-entailment}
    and \eqref{eq:disjunct-into-universal-quantifier-backward-entailment}.
\end{proof}

\begin{thm}
    \label{tautology:grouping-conjuncts-into-universally-quantified-formula}
    Given any well-formed formulas \( \Gamma \) and \( \varphi \)
    with free variables \( \alpha_{1}, \dots, \alpha_{n}\) and
    \(\beta_{1}, \dots, \beta_{m} \) respectively,
    \begin{equation*}
        \Gamma(\alpha_{1}, \dots, \alpha_{n}) \land
        \forall x \big( \varphi(x, \beta_{1}, \dots, \beta_{m}) \big)
        \equiv
        \forall x \big( \Gamma(\alpha_{1}, \dots, \alpha_{n}) \land
        \varphi(x, \beta_{1}, \dots, \beta_{m})\big)
    \end{equation*}
\end{thm}

\begin{proof}
    Let \( \Gamma \) and \( \varphi \) be well-formed formulas interpreted in the context
    of an arbitrary, non-empty domain \( D \). Let \( \alpha_{1}, \dots, \alpha_{n} \) and
    \( \beta_{1}, \dots, \beta_{m} \) represent objects in \( D \).

    \vspace{0.5cm}
    \noindent
    Suppose that \( \Gamma(\alpha_{1}, \dots, \alpha_{n}) \land
    \forall x \big( \varphi(x, \beta_{1}, \dots, \beta_{m}) \big) \) holds.
    Now, let \( x_{0} \in D \) be arbitrary.
    Since \( \forall x \big( \varphi(x, \beta_{1}, \dots, \beta_{m}) \big) \) is true,
    \( \varphi(x_{0}, \beta_{1}, \dots, \beta_{m}) \) follows by Universal Instantiation.
    By Conjunction Introduction and Universal Generalization, it follows that,
    \begin{equation*}
        \forall x \big( \Gamma(\alpha_{1}, \dots, \alpha_{n}) \land
        \varphi(x, \beta_{1}, \dots, \beta_{m})\big)
    \end{equation*}
    Using the Deduction Theorem, we derive the implication:
    \begin{equation*}
        \Gamma(\alpha_{1}, \dots, \alpha_{n}) \land
        \forall x \big( \varphi(x, \beta_{1}, \dots, \beta_{m}) \big)
        \implies
        \forall x \big( \Gamma(\alpha_{1}, \dots, \alpha_{n}) \land
        \varphi(x, \beta_{1}, \dots, \beta_{m})\big)
    \end{equation*}
    Given that the interpretation of all the formulas is arbitrary, it follows that,
    \begin{equation}
        \label{eq:conjunct-into-universal-quantifier-forward-entailment}
        \Gamma(\alpha_{1}, \dots, \alpha_{n})
        \land
        \forall x \big( \varphi(x, \beta_{1}, \dots, \beta_{m}) \big)
        \models
        \forall x \big( \Gamma(\alpha_{1}, \dots, \alpha_{n})
        \land
        \varphi(x, \beta_{1}, \dots, \beta_{m})\big)
    \end{equation}

    \vspace{0.5cm}
    \noindent
    Conversely, if \( \forall x \big( \Gamma(\alpha_{1}, \dots, \alpha_{n}) \land
    \varphi(x, \beta_{1}, \dots, \beta_{m})\big) \) holds, then \( \Gamma(\alpha_{1}, \dots, \alpha_{n}) \)
    must be true. Let
    \( x_{0} \) be an arbitrary object in the domain. Again, \( \varphi(x_{0}, \beta_{1}, \dots, \beta_{m}) \)
    follows directly from the initial hypothesis. Since \( x_{0} \) is arbitrary, Universal Generalization permits
    us to conclude \( \forall x \big( \varphi(x, \beta_{1}, \dots, \beta_{m}) \big) \). By Conjunction Introduction,
    we have,
    \begin{equation*}
        \Gamma(\alpha_{1}, \dots, \alpha_{n}) \land
        \forall x \big( \varphi(x, \beta_{1}, \beta_{m}) \big)
    \end{equation*}
    Using the Deduction Theorem, we infer,
    \begin{equation*}
        \forall x \big( \Gamma(\alpha_{1}, \dots, \alpha_{n}) \land
        \varphi(x, \beta_{1}, \dots, \beta_{m}) \big)
        \implies
        \Gamma(\alpha_{1}, \dots, \alpha_{n})
        \land
        \forall x \big( \varphi(x, \beta_{1}, \dots, \beta_{m}) \big)
    \end{equation*}
    And since the interpretation is arbitrary,
    \begin{equation}
        \label{eq:conjunct-into-universal-quantifier-backward-entailment}
        \forall x \big( \Gamma(\alpha_{1}, \dots, \alpha_{n}) \land
        \varphi(x, \beta_{1}, \dots, \beta_{m}) \big)
        \models
        \Gamma(\alpha_{1}, \dots, \alpha_{n})
        \land
        \forall x \big( \varphi(x, \beta_{1}, \dots, \beta_{m}) \big)
    \end{equation}
    Together, \eqref{eq:conjunct-into-universal-quantifier-forward-entailment} and \eqref{eq:conjunct-into-universal-quantifier-backward-entailment}
    establish the logical equivalence.
\end{proof}

\begin{thm}
    \label{tautology:grouping-disjuncts-into-existentially-quantified-formula}
    For any well-formed formula \( \Gamma \) and \( \varphi \) with free variables
    \( \alpha_{1}, \dots \alpha_{n} \) and \( \beta_{1}, \dots, \beta_{m} \) respectively,
    \begin{equation*}
        \Gamma(\alpha_{1}, \dots, \alpha_{n})
        \lor
        \exists x \big( \varphi(x, \beta_{1}, \dots, \beta_{m}) \big)
        \equiv
        \exists x \big( \Gamma(\alpha_{1}, \dots, \alpha_{n}) \lor
        \varphi(x, \beta_{1}, \dots, \beta_{m}) \big)
    \end{equation*}
\end{thm}

\begin{proof}
    Let \( \Gamma \) and \( \varphi \) be well-formed formulas interpreted in the context of an arbitrary, non-empty domain \( D \).
    Let \( \alpha_{1}, \dots, \alpha_{n} \) and \( \beta_{1}, \dots, \beta_{m} \) represent objects in \( D \).

    \vspace{0.5cm}
    \noindent
    Suppose that \( \Gamma(\alpha_{1}, \dots, \alpha_{n}) \lor \exists x \big( \varphi(x, \beta_{1}, \dots, \beta_{m}) \big) \)
    is true, which is only possible in three scenarios.

    \vspace{0.5cm}
    \noindent
    {\scshape Case 1}: \( \Gamma(\alpha_{1}, \dots, \alpha_{n}) \land \neg \, \exists x \big( \varphi(x, \beta_{1}, \dots, \beta_{m}) \big) \).

    \vspace{0.1cm}
    \noindent
    Since \( \Gamma(\alpha_{1}, \dots, \alpha_{n}) \) is true, it follows that \( \forall x \big( \Gamma(\alpha_{1}, \dots, \alpha_{n})
    \lor \varphi(x, \beta_{1}, \dots, \beta_{m}) \big) \) is true. Since the domain \( D \) is non-empty,
    \( \exists x \big( \Gamma(\alpha_{1}, \dots, \alpha_{n}) \lor
    \varphi(x, \beta_{1}, \dots, \beta_{m}) \big) \) is necessarily true.

    \vspace{0.5cm}
    \noindent
    {\scshape Case 2}: \( \neg \, \Gamma(\alpha_{1}, \dots, \alpha_{n}) \land
    \exists x \big( \varphi(x, \beta_{1}, \dots, \beta_{m}) \big) \)

    \vspace{0.1cm}
    \noindent
    Let \( x_{0} \) be an object in the domain such that \( \varphi(x_{0}, \beta_{1}, \dots, \beta_{m}) \). By Disjunction Introduction,
    clearly \( \Gamma(\alpha_{1}, \dots, \alpha_{n}) \lor \varphi(x_{0}, \beta_{1}, \dots, \beta_{m}) \) is true. And by Existential
    Introduction, we have, \( \exists x \big( \Gamma(\alpha_{1}, \dots, \alpha_{n}) \lor \varphi(x, \beta_{1}, \dots, \beta_{m}) \big) \).

    \vspace{0.5cm}
    \noindent
    {\scshape Case 3}: \( \Gamma(\alpha_{1}, \dots, \alpha_{n}) \land \exists x \big( \varphi(x, \beta_{1}, \dots, \beta_{m}) \big) \)

    \vspace{0.1cm}
    \noindent
    Since \( \Gamma(\alpha_{1}, \dots, \alpha_{n}) \) is true, this reduces to the first case.

    \vspace{0.5cm}
    \noindent
    Evidently, \( \exists x \big( \Gamma(\alpha_{1}, \dots, \alpha_{n}) \lor \varphi(x, \beta_{1}, \dots, \beta_{m}) \big) \) is implied
    in all three cases. Using the Deduction Theorem, we conclude,
    \begin{equation*}
        \Gamma(\alpha_{1}, \dots, \alpha_{n}) \lor
        \exists x \big( \varphi(x, \beta_{1}, \dots, \beta_{m}) \big)
        \implies
        \exists x \big( \Gamma(\alpha_{1}, \dots, \alpha_{n})
        \lor
        \varphi(x, \beta_{1}, \dots, \beta_{m}) \big)
    \end{equation*}
    Since the implication was derived using an arbitrary interpretation, we can derive the following entailment:
    \begin{equation}
        \label{eq:disjunct-into-existential-quantifier-forward-entailment}
        \Gamma(\alpha_{1}, \dots, \alpha_{n}) \lor
        \exists x \big( \varphi(x, \beta_{1}, \dots, \beta_{m}) \big)
        \models
        \exists x \big( \Gamma(\alpha_{1}, \dots, \alpha_{n}) \lor
        \varphi(x, \beta_{1}, \dots, \beta_{m}) \big)
    \end{equation}

    \vspace{0.5cm}
    \noindent
    Proving the converse is also fairly straightforward. Naturally, if
    \begin{equation*}
        \exists x \big( \Gamma(\alpha_{1}, \dots, \alpha_{n}) \lor \varphi(x, \beta_{1}, \dots, \beta_{m}) \big)
    \end{equation*}
    is true, then there are only two possible scenarios to consider\,\footnote{Again, this is the Law of Excluded Middle at play.}.

    \vspace{0.5cm}
    \noindent
    {\scshape Case 1}: \( \Gamma(\alpha_{1}, \dots, \alpha_{n}) \)

    \vspace{0.1cm}
    \noindent
    {\scshape Case 2}: \( \neg \, \Gamma(\alpha_{1}, \dots, \alpha_{n}) \)

    \vspace{0.5cm}
    \noindent
    In the first case, \( \Gamma(\alpha_{1}, \dots, \alpha_{n}) \lor \exists x \big( \varphi(x, \beta_{1}, \dots, \beta_{m}) \big) \) follows
    readily from Disjunction Introduction.
    In the second case, we use Disjunctive Syllogism to conclude that,
    \begin{equation*}
        \exists x \big[ \varphi(x, \beta_{1}, \dots, \beta_{m}) \big]
    \end{equation*}
    Again, \( \Gamma(\alpha_{1}, \dots, \alpha_{n}) \lor \exists x \mleft( \varphi(x, \beta_{1}, \dots, \beta_{m}) \mright) \) follows from Disjunction
    Introduction. Applying the Deduction Theorem, we infer,
    \begin{equation*}
        \exists x \big( \Gamma(\alpha_{1}, \dots, \alpha_{n}) \lor \varphi(x, \beta_{1}, \dots, \beta_{m}) \big)
        \implies
        \Gamma(\alpha_{1}, \dots, \alpha_{n}) \lor
        \exists x \big( \varphi(x, \beta_{1}, \dots, \beta_{m}) \big)
    \end{equation*}
    And since the implication is independent of implication, we have,
    \begin{equation}
        \label{eq:disjunct-into-existential-quantifier-backward-implication}
        \exists x \big( \Gamma(\alpha_{1}, \dots, \alpha_{n}) \lor \varphi(x, \beta_{1}, \dots, \beta_{m}) \big)
        \models
        \Gamma(\alpha_{1}, \dots, \alpha_{n}) \lor
        \exists x \big( \varphi(x, \beta_{1}, \dots, \beta_{m}) \big)
    \end{equation}
    Together, \eqref{eq:disjunct-into-existential-quantifier-forward-entailment} and \eqref{eq:disjunct-into-existential-quantifier-backward-implication}
    establish the logical equivalence.
\end{proof}

\begin{thm}
    \label{tautology:grouping-conjuncts-into-existentially-quantified-formula}
    For any well-formed formula \( \Gamma \) and \( \varphi \) with free variables
    \( \alpha_{1}, \dots, \alpha_{n} \) and \( \beta_{1}, \dots, \beta_{m} \) respectively,
    \begin{equation*}
        \Gamma(\alpha_{1}, \dots, \alpha_{n}) \land
        \exists x \big( \varphi(x, \beta_{1}, \dots, \beta_{m}) \big)
        \equiv
        \exists x \big( \Gamma(\alpha_{1}, \dots, \alpha_{n}) \land
        \varphi(x, \beta_{1}, \dots, \beta_{m}) \big)
    \end{equation*}
\end{thm}

\begin{proof}
    Let \( \Gamma \) and \( \varphi \) be well-formed fomulas interpreted in the context of an
    arbitrary, non-empty domain, \( D \). Let \( \alpha_{1}, \dots, \alpha_{n} \) and
    \( \beta_{1}, \dots, \beta_{m} \) represent objects in \( D \).

    \vspace{0.5cm}
    \noindent
    Suppose that \( \Gamma(\alpha_{1}, \dots, \alpha_{n}) \land \exists x \mleft( \varphi(x, \beta_{1}, \dots, \beta_{m}) \mright) \)
    is true. Clearly, \( \varphi(x_{0}, \beta_{1}, \dots, \beta_{m}) \) is true for some \( x_{0} \in D \). Thus, by Conjunction Introduction,
    \begin{equation*}
        \Gamma(\alpha_{1}, \dots, \alpha_{n}) \land \varphi(x_{0}, \beta_{1}, \dots, \beta_{m})
    \end{equation*}
    And by Existential Instantiation,
    \begin{equation*}
        \exists x \big( \Gamma(\alpha_{1}, \dots, \alpha_{n}) \land \varphi(x, \beta_{1}, \dots, \beta_{m}) \big)
    \end{equation*}
    Thus,
    \begin{equation*}
        \Gamma(\alpha_{1}, \dots, \alpha_{n}) \land
        \exists x \big( \varphi(x, \beta_{1}, \dots, \beta_{m}) \big)
        \models
        \exists x \big( \Gamma(\alpha_{1}, \dots, \alpha_{n}) \land
        \varphi(x, \beta_{1}, \dots, \beta_{m})\big)
    \end{equation*}

    \vspace{0.5cm}
    \noindent
    Proving the converse is trivial. If \( \exists x \big( \Gamma(\alpha_{1}, \dots, \alpha_{n}) \land
    \varphi(x, \beta_{1}, \dots, \beta_{m})\big) \) is true, then clearly,
    \begin{equation*}
        \Gamma(\alpha_{1}, \dots, \alpha_{n}) \land \varphi(x_{0}, \beta_{1}, \dots, \beta_{m})
    \end{equation*}
    for some \( x_{0} \in D \). By Conjunction Elimination, obviously \( \varphi(x_{0}, \beta_{1}, \dots, \beta_{m}) \)
    is true for some \( x_{0} \in D \). By Existential Generalization, we thus conclude,
    \begin{equation*}
        \Gamma(\alpha_{1}, \dots, \alpha_{n}) \land
        \exists x \big( \varphi(x, \beta_{1}, \dots, \beta_{m}) \big)
    \end{equation*}
    And therefore,
    \begin{equation*}
        \exists x \big( \Gamma(\alpha_{1}, \dots, \alpha_{n}) \land
        \varphi(x, \beta_{1}, \dots, \beta_{m})\big)
        \models
        \Gamma(\alpha_{1}, \dots, \alpha_{n}) \land
        \exists x \big( \varphi(x, \beta_{1}, \dots, \beta_{m}) \big) \qedhere
    \end{equation*}
\end{proof}

\begin{thm}
    \label{tautology:combining-unquantified-antecedent-and-universally-quantified-conseqeunt-into-a-single-quantified-formula}
    For any well-formed formula \( \Gamma \) and \( \varphi \) with free variables
    \( \alpha_{1}, \dots, \alpha_{n} \) and \( \beta_{1}, \dots, \beta_{m} \) respectively,
    \begin{equation*}
        \Gamma(\alpha_{1}, \dots, \alpha_{n})
        \implies
        \forall x \big[ \varphi(x, \beta_{1}, \dots, \beta_{m}) \big]
        \equiv
        \forall x \big[ \Gamma(\alpha_{1}, \dots, \alpha_{n}) \implies \varphi(x, \beta_{1}, \dots, \beta_{m}) \big]
    \end{equation*}
\end{thm}

\begin{proof}
    Let \( \Gamma \) and \( \varphi \) be well-formed formulas interpreted in the context of an arbitrary, non-empty
    domain \( D \). Let \( \alpha_{1}, \dots, \alpha_{n} \) and \( \beta_{1}, \dots, \beta_{m} \) be objects in \( D \).
    Using the truth table for logical implication and \cref{tautology:grouping-disjuncts-into-universally-quantified-formula},
    we transform one of the statement into the other like so:
    \begin{align*}
        \Gamma(\alpha_{1}, \dots, \alpha_{n})
        {\implies}
        \forall x \big( \varphi(x, \beta_{1}, \dots, \beta_{m}) \big)
        &\iff
        \neg \, \Gamma(\alpha_{1}, \dots, \alpha_{n})
        \lor
        \forall x \big( \varphi(x, \beta_{1}, \dots, \beta_{m}) \big)\\
        &\iff
        \forall x \big( \neg \, \Gamma(\alpha_{1}, \dots, \alpha_{n})
        \lor
        \varphi(x, \beta_{1}, \dots, \beta_{m}) \big)\\
        &\iff
        \forall x \big( \Gamma(\alpha_{1}, \dots, \alpha_{n})
        {\implies}
        \varphi(x, \beta_{1}, \dots, \beta_{m}) \big)
    \end{align*}
    The biconditional holds regardless of the interpretation of choice. Thus, the two statements are logically equivalent.
\end{proof}

\begin{thm}
    \label{tautology:combining-universally-quantified-antecedent-and-unquantified-consequent-into-a-single-quantified-formula}
    For any well-formed formula \( \Gamma \) and \( \varphi \) with free variables
    \( \alpha_{1}, \dots, \alpha_{n} \) and \( \beta_{1}, \dots, \beta_{m} \) respectively,
    \begin{equation*}
        \forall x \big( \varphi(x, \beta_{1}, \dots, \beta_{m}) \big)
        {\implies}
        \Gamma(\alpha_{1}, \dots, \alpha_{n})
        \equiv
        \exists x \big( \varphi(x, \beta_{1}, \dots, \beta_{m})
        {\implies}
        \Gamma(\alpha_{1}, \dots, \alpha_{n})
        \big)
    \end{equation*}
\end{thm}

\begin{proof}
    Let \( \Gamma \) and \( \varphi \) be well-formed formulas interpreted in the context of an arbitrary,
    non-empty domain \( D \). Let \( \alpha_{1}, \dots, \alpha_{n} \) and \( \beta_{1}, \dots, \beta_{m} \)
    be objects in \( D \). Using the truth table of the logical implication and
    \cref{tautology:grouping-disjuncts-into-universally-quantified-formula}, we can transform one statement into
    another like so:
    \begin{align*}
        \forall x \big( \varphi(x, \beta_{1}, \dots, \beta_{m}) \big)
        {\implies}
        \Gamma(\alpha_{1}, \dots, \alpha_{n})
        &\iff
        \neg \, \forall x \big( \varphi(x, \beta_{1}, \dots, \beta_{m}) \big)
        \lor
        \Gamma(\alpha_{1}, \dots, \alpha_{n})\\
        &\iff
        \exists x \big( \neg \, \varphi(x, \beta_{1}, \dots, \beta_{m}) \big)
        \lor
        \Gamma(\alpha_{1}, \dots, \alpha_{n})\\
        &\iff
        \exists x \big( \neg \, \varphi(x, \beta_{1}, \dots, \beta_{m})
        \lor
        \Gamma(\alpha_{1}, \dots, \alpha_{n}) \big)\\
        &\iff
        \exists x \big( \varphi(x, \beta_{1}, \dots, \beta_{m})
        {\implies}
        \Gamma(\alpha_{1}, \dots, \alpha_{n})\big)
    \end{align*}
    Clearly, the biconditional holds for any interpretation. Therefore, the statements are indeed
    logically equivalent.
\end{proof}

\begin{thm}
    \label{tautology:using-a-universally-quantified-formula-to-express-the-fact-that-an-object-satisfies-a-formula}
    For any well-formed formula \( \varphi \) and free variables
    \( \alpha, \beta_{1}, \dots, \beta_{m} \),
    \begin{equation*}
        \forall x \big( x = \alpha \implies \varphi(x, \beta_{1}, \dots, \beta_{m}) \big)
        \equiv
        \varphi(\alpha, \beta_{1}, \dots, \beta_{m})
    \end{equation*}
\end{thm}

\begin{proof}
    Let \( \varphi \) be a well-formed formula interpreted in the context of a non-empty
    domain, \( D \). Let \( \alpha, \beta_{1}, \dots, \beta_{m} \) represent objects in \( D \).
    Suppose that the well-formed formula
    \( \forall x \big( x = \alpha \implies \varphi(x, \beta_{1}, \dots, \beta_{m}) \big) \) is true.
    Then by Universal Instantiation, we have,
    \begin{equation*}
        \alpha = \alpha \implies \varphi(\alpha, \beta_{1}, \dots, \beta_{m})
    \end{equation*}
    Since \( \alpha = \alpha \) is trivially true, we conclude by Modus Ponens that
    \( \varphi(\alpha, \beta_{1}, \dots, \beta_{m}) \) is also true.
    \begin{equation*}
        \therefore \, \forall x \big( x = \alpha \implies \varphi(x, \beta_{1}, \dots, \beta_{m}) \big)
        \implies
        \varphi(\alpha, \beta_{1}, \dots, \beta_{m})
    \end{equation*}
    Since the domain set \( D \) (the interpretation) is arbitrary selected, it follows that,
    \begin{equation}
        \label{eq:all-x-equals-alpha-implies-phi-x-entails-phi-alpha}
        \forall x \big( x = \alpha \implies \varphi(x, \beta_{1}, \dots, \beta_{m}) \big)
        \models
        \varphi(\alpha, \beta_{1}, \dots, \beta_{m})
    \end{equation}
    Conversely, suppose that \( \varphi(\alpha, \beta_{1}, \dots, \beta_{m}) \) is true for an arbitrary interpretation.
    Let \( x_{0} \in D \) such that \( x_{0} = \alpha \).
    Since \( \varphi(\alpha, \beta_{1}, \dots, \beta_{m}) \) is true and \( x_{0} = \alpha \) ,
    it follows that \( \varphi(x_{0}, \beta_{1}, \dots, \beta_{m}) \) is also true.
    Thus, \( x_{0} = \alpha \implies \varphi(x_{0}, \beta_{1}, \dots, \beta_{m})\).
    Since \( x_{0} \) is arbitrary, Universal Generalization yields
    \( \forall x \big( x = \alpha \implies \varphi(x, \beta_{1}, \dots, \beta_{m}) \big) \).
    \begin{equation*}
        \therefore \, \varphi(\alpha, \beta_{1}, \dots, \beta_{m})
        \implies
        \forall x \big( x = \alpha \implies \varphi(x, \beta_{1}, \dots, \beta_{m}) \big)
    \end{equation*}
    Again, since the interpretation is arbitrary, it follows that,
    \begin{equation}
        \label{eq:phi-alpha-entails-all-x-equals-alpha-implies-phi-x}
        \varphi(\alpha, \beta_{1}, \dots, \beta_{m})
        \models
        \forall x \big( x = \alpha \implies \varphi(x, \beta_{1}, \dots, \beta_{m}) \big)
    \end{equation}
    Together, \eqref{eq:all-x-equals-alpha-implies-phi-x-entails-phi-alpha} and \eqref{eq:phi-alpha-entails-all-x-equals-alpha-implies-phi-x}
    establish the equivalence as needed.
\end{proof}

\begin{thm}
    \label{tautology:using-an-existentially-quantified-formula-to-express-the-fact-that-an-object-satisfies-a-formula}
    For any well-formed formula \( \varphi \) and free variables \( \alpha, \beta_{1}, \dots, \beta_{m} \),
    \begin{equation*}
        \exists x \big( x = \alpha \land \varphi(x, \beta_{1}, \dots, \beta_{m}) \big)
        \equiv
        \varphi(\alpha, \beta_{1}, \dots, \beta_{m})
    \end{equation*}
\end{thm}

\begin{proof}
    Let \( \varphi \) be a well-formed formula interpreted in the context of a non-empty domain, \( D \).
    Let \( \alpha, \beta_{1}, \dots, \beta_{m} \) represent objects in \( D \). Supposed that the well-formed
    formula \( \exists x \mleft( x = \alpha \land \varphi(x, \beta_{1}, \dots, \beta_{m}) \mright) \)
    is true. Then clearly,
    \begin{equation*}
        x_{0} = \alpha \land \varphi(x_{0}, \beta_{1}, \dots, \beta_{m})
    \end{equation*}
    for some \( x_{0} \in D \). Since \( x_{0} = \alpha \), it follows that
    \( \varphi(\alpha, \beta_{1}, \dots, \beta_{m}) \) must also be true.
    In other words,
    \begin{equation*}
        \exists x \big( x = \alpha \land \varphi(x, \beta_{1}, \dots, \beta_{m}) \big)
        \models
        \varphi(\alpha, \beta_{1}, \dots, \beta_{m})
    \end{equation*}
    Conversely, if \( \varphi(\alpha, \beta_{1}, \dots, \beta_{m}) \) is true, then clearly
    \( \alpha = \alpha \land \varphi(\alpha, \beta_{1}, \dots, \beta_{m}) \). Since
    \( \alpha \in D \), then we can use Existential Generalization to conclude,
    \begin{equation*}
        \exists x \big( x = \alpha \land \varphi(x, \beta_{1}, \dots, \beta_{m}) \big)
    \end{equation*}
    Thus,
    \begin{equation*}
        \varphi(\alpha, \beta_{1}, \dots, \beta_{m})
        \models
        \exists x \big( x = \alpha \land \varphi(x, \beta_{1}, \dots, \beta_{m}) \big)
    \end{equation*}
    Since the entailment has been established in both directions, the two formulas are indeed logically
    equivalent.
\end{proof}
